% !TeX root = ../main.tex
% Este capítulo se puede dividir en varios (diseño hardware, software,
% modelo…). Crea un fichero por capítulo y añádelo con \include en main.tex.
\chapter{Desarrollo}\label{cap:desarrollo}

Describe la implementación con el detalle suficiente para que otra
persona pueda reproducirla, y justifica las decisiones de diseño.

\section{Algoritmos}

\begin{algorithm}[H]
  \caption{Media móvil exponencial.}
  \label{alg:ema}
  \KwIn{señal $x_1,\dots,x_N$, factor $\alpha \in (0,1]$}
  \KwOut{señal filtrada $y_1,\dots,y_N$}
  $y_1 \gets x_1$\;
  \For{$k \gets 2$ \KwTo $N$}{
    $y_k \gets \alpha x_k + (1-\alpha)\, y_{k-1}$\;
  }
  \Return{$y$}
\end{algorithm}

\section{Código}

Incluye solo los fragmentos imprescindibles (\cref{lst:ema}). El código
completo debe estar en el repositorio y se referencia desde el
\cref{anx:codigo}.

\begin{lstlisting}[language=Python, caption={Implementación del \cref{alg:ema}.}, label={lst:ema}]
def ema(x, alpha):
    """Media móvil exponencial."""
    y = [x[0]]
    for xk in x[1:]:
        y.append(alpha * xk + (1 - alpha) * y[-1])
    return y
\end{lstlisting}
